% Study sheet: Volume III, Chapter 19, Agent Economics.
% Numbers: \val{...} from agent_economics_studysheet_numbers.tex, written by agent_economics_studysheet_numbers.py.
\documentclass{studysheet}
\input{agent_economics_studysheet_numbers}
\sheetvolume{3}
\sheetchapter{19}{Agent Economics}
\sheetsources{Numbers from the chapter's mlsysim scenarios (hypothetical prices, cohorts and timings); questions marked assumed use new inputs.}

\begin{document}
% ===================================================================== FRONT: LEARN
\frontside
\begin{bigidea}
An agent service delivers \hl{accepted tasks}, not tokens. Every attempt, failed or canceled,
is charged to the tasks one declared contract accepts, so \hl{acceptance divides the whole bill} while a
token discount trims only the model's share: buy a lever only if it lowers \hl{cost per accepted task}
or meets a deadline the objective values.
\end{bigidea}

\begin{leftcol}
\sheetsec{Key ideas}
\begin{ideas}
\idea{Acceptance divides the whole bill.}{A token discount trims one term; acceptance divides them all:
  sandbox, tools, checks, review.}
\idea{The cohort is the account.}{Every initiated task, failed, canceled or pending, stays on the bill;
  retries add no acceptances.}
\idea{Caching lowers the coefficient, not the growth.}{Prefix reuse bills the quadratic re-send at the
  cached rate; compaction pays only with enough calls left.}
\idea{A cascade is only as cheap as its gates.}{Every reached tier pays its call and its gate; savings come
  from early exits, false accepts priced in.}
\idea{Utilization is paid in waiting.}{Wait climbs steeply near saturation; a trajectory pays it on every
  serial call.}
\idea{Reserve before dispatch, release at settlement.}{Admission atomically reserves a charge bound; a
  stopped child releases allocation minus settled minus that bound.}
\end{ideas}

\vspace{8pt}
\fcard{Cohort cost per accepted task}%
  {\CostPerAcceptedTask = \frac{\sum_{i=1}^{\InitiatedTasks} {\TaskCost}_i}{\AcceptedTasks}
   = \AcceptedTaskMeanCost + \frac{1-\AcceptanceFraction}{\AcceptanceFraction}\,\UnacceptedTaskMeanCost}%
  {$\InitiatedTasks$ tasks initiated, $\AcceptedTasks$ accepted, $\AcceptanceFraction = \AcceptedTasks/\InitiatedTasks$;
   $\AcceptedTaskMeanCost$ and $\UnacceptedTaskMeanCost$ are the mean costs of accepted and unaccepted
   tasks.}%
  {}{}
\fcard{Compaction payback}%
  {\CompactionUpfrontCost = \OutputTokenPrice\SummaryOutputLength + \InputTokenPrice\CompactedContextLength,
   \qquad \CompactionPerCallSaving = \CachedInputTokenPrice(\SequenceLength-\CompactedContextLength)}%
  {Compacting $\SequenceLength$ tokens to $\CompactedContextLength$: upfront cost $\CompactionUpfrontCost$ (a
   $\SummaryOutputLength$-token summary, one uncached read), saving $\CompactionPerCallSaving$ per later call. With both
   positive, it pays once $1+\lceil\CompactionUpfrontCost/\CompactionPerCallSaving\rceil$ calls remain, first compacted call included.}%
  {Per-token prices: fresh $\InputTokenPrice$, cached $\CachedInputTokenPrice$, output $\OutputTokenPrice$.}{}
\end{leftcol}%
\begin{rightcol}
\fcard{Speed pays only on the critical path}%
  {\AcceleratedEpisodeTime = \Time\left[(1-\ModelTimeFraction) + \frac{\ModelTimeFraction}{\AcceleratorGain}\right],
   \qquad \text{ceiling } \frac{1}{1-\ModelTimeFraction}}%
  {Model calls take fraction $\ModelTimeFraction$ of the critical path $\Time$, summed as disjoint intervals;
   a $\AcceleratorGain$-fold faster model leaves the rest untouched.}%
  {}%
  {\begin{tikzpicture}[sheet]
     \pgfmathsetmacro\sc{176/(\val{barTests}+\val{barSandbox}+\val{barModel}+\val{barRuntime})}
     \pgfmathsetmacro\xa{\val{barTests}*\sc}
     \pgfmathsetmacro\xb{\xa+\val{barSandbox}*\sc}
     \pgfmathsetmacro\xc{\xb+\val{barModel}*\sc}
     \pgfmathsetmacro\xd{\xb+\val{barModel}*\sc/\val{cpgain}}
     \draw[nbox] (0,22) rectangle (\xa,34);
     \draw[accbox] (\xa,22) rectangle (\xb,34);
     \draw[hotbox] (\xb,22) rectangle (\xc,34);
     \draw[nbox] (0,0) rectangle (\xa,12);
     \draw[accbox] (\xa,0) rectangle (\xb,12);
     \draw[hotbox] (\xb,0) rectangle (\xd,12);
     \node[lab] at ({\xa/2},28) {tests \val{cptests} s};
     \node[lab] at ({\xa/2},6) {tests};
     \draw[dim] ({(\xa+\xb)/2},34.5) -- ({(\xa+\xb)/2-7},39);
     \draw[dim] ({(\xb+\xc)/2},34.5) -- ({(\xb+\xc)/2+7},39);
     \node[small, anchor=south east, inner sep=1pt] at ({(\xa+\xb)/2-7},39) {sandbox};
     \node[small, anchor=south west, inner sep=1pt] at ({(\xb+\xc)/2+7},39) {model};
     \node[note, anchor=west] at ({\xc+3},28) {\val{cptotal} s};
     \node[note, anchor=west] at ({\xd+3},6) {\val{cpnew} s};
     \node[note, anchor=north west] at (0,-3) {model \val{cpmodel} s (\val{cpfrac}\%): \val{cpgain}$\times$ faster
       gives \val{cpspeed}$\times$; instant, \val{cpcap}$\times$};
   \end{tikzpicture}}
\fcard{Goodput sets the load; utilization sets the wait}%
  {\begin{gathered}\TrajectoryAdmissionRate = \frac{\TrajectoryGoodput}{\AcceptanceFraction}, \qquad
   \ConcurrentTrajectories = \TrajectoryAdmissionRate\,\mathbb{E}[\TrajectoryDuration]\\[2pt]
   \ArrivalRate = \TrajectoryAdmissionRate\,\mathbb{E}[\TrajectoryCallCount], \qquad
   \ServingUtilization = \ArrivalRate\,\mathbb{E}[\RequestServiceTime]\\[3pt]
   \mathbb{E}[\RequestQueueTime] \approx \frac{\InterarrivalVariationCoefficient^2+\ServiceTimeVariationCoefficient^2}{2}\,
   \frac{\ServingUtilization}{1-\ServingUtilization}\,\mathbb{E}[\RequestServiceTime]\end{gathered}}%
  {For goodput $\TrajectoryGoodput$ (accepted tasks per minute), admit $\TrajectoryAdmissionRate$; Little's law
   gives $\ConcurrentTrajectories$ present; calls arrive at $\ArrivalRate$. At one server, burstiness
   ($\InterarrivalVariationCoefficient^2$, $\ServiceTimeVariationCoefficient^2$) and service time
   $\RequestServiceTime$ scale the wait.}%
  {}{}
\fcard{Loop cap with a human fallback}%
  {\text{run turn } \TurnIndex \text{ only if } \StageConditionalAcceptanceProbability_{\TurnIndex}\,\HumanFallbackCost > \UniformLoopTurnCost}%
  {$\StageConditionalAcceptanceProbability_{\TurnIndex}$: chance turn $\TurnIndex$ is accepted, given it is
   reached. $\HumanFallbackCost$: a fallback that resolves every leftover ticket. $\UniformLoopTurnCost$: the
   full turn, check and sandbox included.}%
  {}{}
\end{rightcol}

% ===================================================================== BACK: PRACTICE
\backside
\begin{answers}
\ans{Finish}{\$\val{bcohort} $\div$ \val{fixaccepted} = \$\val{fixceff}, still above \$\val{aceff}. Beating it
  needs more than \val{bcohort} $\div$ \val{aceff} $\approx$ \val{needacc}: \val{needwhole} (cohort card).}
\ans{Critical path}{(a) $\ModelTimeFraction \approx$ \val{qfrac}\%; \val{qrest} + \val{qmodel} $\div$ \val{qgain} =
  \val{qnew} s, \val{qspeed}$\times$. (b) \val{qtotal} $\div$ \val{qrest} = \val{qcap}$\times$. (c) Halving tests:
  \val{qtestsave} s against \val{qmodelsave} s (critical-path card).}
\ans{Capacity}{(a) \val{sgood} $\div$ \val{salpha} = \val{sadmit}; $\times$ \val{sstay} = \val{sconc}
  present; $\times$ \val{scalls} = \val{scallrate} calls per minute. (b) $\ServingUtilization$ = \val{scallsec} per s $\times$ \val{ssvc} =
  \val{srho}; \val{svar} $\times$ \val{sratio} $\times$ \val{ssvc} = \val{swait} s. (c) $\ServingUtilization$ = \val{srho} $\times$ \val{salpha} $\div$ \val{salphatwo}
  $\approx$ \val{srhotwo}; wait factor \val{sratio} becomes \val{sratiotwo}: \val{swaittwo} s,
  \val{swaitx}$\times$ for \val{sloadrise}\% more load (capacity card).}
\columnbreak
\ans{Loop cap}{(a) $\StageConditionalAcceptanceProbability_{\TurnIndex}\HumanFallbackCost$ = \val{lgone},
  \val{lgtwo}, \val{lgthree}, \val{lgfour} against \$\val{lturn}: cap at \val{lcapword} turns. (b) \val{lhone},
  \val{lhtwo}, \val{lhthree}: cap at \val{lcaptwoword} (loop card).}
\ans{Budget}{(a) \val{balloc} $-$ \val{bsettled} $-$ \val{bbound} = \$\val{bsafe}, not \$\val{bnaive}. (b) \val{bbound} $-$ \val{blate} = \$\val{bmore} more; \$\val{bclose} in all. (reserve idea).}
\end{answers}

\begin{leftcol}
\sheetsec{Worked example}
\begin{practicebox}\boxtitle{Does a tenth-price model make repairs cheaper?}
\val{tickets} tickets, one attempt each, one contract. \textbf{Frontier:} \val{aturns} turns, tokens
\$\val{amodel} per attempt, \val{aaccepted} accepted. \textbf{Low price} (token rates about a tenth): \val{bturns}
turns, tokens \$\val{bmodel}, \val{baccepted} accepted. Both pay \$\val{envturn} environment per turn, \$\val{assess}
assessment per attempt.\par\vspace{5pt}
{\centering
\begin{tikzpicture}[sheet]
  \pgfmathsetmacro\sc{150/(\val{barBmodel}+\val{barBenv}+\val{barBassess})}
  \foreach \lab/\m/\e/\c/\y/\tot in {frontier/\val{barAmodel}/\val{barAenv}/\val{barAassess}/18/\val{aceff},
                                     low price/\val{barBmodel}/\val{barBenv}/\val{barBassess}/0/\val{bceff}}{
    \pgfmathsetmacro\xa{48+\m*\sc}
    \pgfmathsetmacro\xb{\xa+\e*\sc}
    \pgfmathsetmacro\xc{\xb+\c*\sc}
    \node[lab, anchor=east] at (45,\y+6) {\lab};
    \draw[accbox] (48,\y) rectangle (\xa,\y+12);
    \draw[hotbox] (\xa,\y) rectangle (\xb,\y+12);
    \draw[nbox] (\xb,\y) rectangle (\xc,\y+12);
    \node[note, anchor=west] at ({\xc+3},\y+6) {\$\tot};}
  \pgfmathsetmacro\ma{48+\val{barBmodel}*\sc}
  \pgfmathsetmacro\mb{\ma+\val{barBenv}*\sc}
  \node[small, anchor=north] at ({(48+\ma)/2},-2) {model};
  \node[small, anchor=north] at ({(\ma+\mb)/2},-2) {environment};
  \node[small, anchor=north] at ({\mb+11},-2) {check};
\end{tikzpicture}\par}\vspace{2pt}
{\centering\fontsize{8}{10}\selectfont\color{muted}cost per accepted repair, by term\par}\vspace{3pt}
\begin{steps}
\item Frontier: \val{amodel} + \val{aturns} $\times$ \val{envturn} + \val{assess} = \$\val{aattempt};
  $\times$ \val{tickets} $\div$ \val{aaccepted} = \res{\$\val{aceff}}.
\item Low price: \val{bmodel} + \val{bturns} $\times$ \val{envturn} + \val{assess} = \$\val{battempt};
  $\times$ \val{tickets} $\div$ \val{baccepted} = \res{\$\val{bceff}}.
\item \val{bceff} $\div$ \val{aceff} = \res{\val{ceffratio}$\times$}. Token spend per attempt fell
  \val{modelcut}\%, attempt cost \val{attemptcut}\%, acceptance \val{acceptcut}\%.
\end{steps}
\takeaway{A discount cuts one term; lost acceptance multiplies all.}
\end{practicebox}

\vspace{10pt}
\sheetsec{Your turn: finish it}
\begin{practicebox}\boxtitle{Does a tool fix rescue the low-price model?}
\assumed{} a tool fix lifts the low-price configuration to \val{fixaccepted} accepted, all else
equal.\par\vspace{4pt}
\textbf{Given:} the cohort still costs \val{tickets} $\times$ \$\val{battempt} = \$\val{bcohort}.\par
\textbf{You:} find its cost per accepted repair, and the fewest acceptances that beat \$\val{aceff}.\par\vspace{30pt}
\end{practicebox}
\end{leftcol}%
\begin{rightcol}
\sheetsec{Common mistakes}
\mistake{``Compaction shrinks every later call, so it always saves.''}%
  {At \$\val{cpin} fresh, \$\val{cpcache} cached, \$\val{cpout} output per million tokens, compacting \val{cpctx} tokens
   to \val{cpto} with a \val{cpsum}-token summary: $\CompactionUpfrontCost$ = \$\val{cpup}, $\CompactionPerCallSaving$ = \$\val{cpsave};
   payback $1+\lceil\val{cpratio}\rceil$ = \val{cppay} calls, inside the \val{cprest} left. Halve the cached rate: \val{cppayhalf}, a loss.}{}
\mistake{``A cheap first tier makes the cascade's saving safe.''}%
  {A cascade costs \$\val{cascceff} per accepted task against \$\val{directceff} direct. False accepts on
   \val{fapct}\% of its \val{tieronexits} first-tier exits, at \$\val{harmeach} each, add \$\val{harmtotal}
   over \val{cascacc} accepted: \$\val{cascharm}, dearer than direct.}{}

\vspace{2pt}
\sheetsec{Check yourself}
\begin{qs}
\item \assumed{} a serial critical path: \val{qmodel} s model, \val{qtests} s tests, \val{qsandbox} s sandbox.
  (a) Time and speedup, model \val{qgain}$\times$ faster? (b) The most any model speedup buys?
  (c) Which saves more: the \val{qgain}$\times$ model or halved tests?
\item \assumed{} a support service delivers \val{sgood} accepted sessions per minute at
  $\AcceptanceFraction = \val{salpha}$; a session stays \val{sstay} minutes and makes \val{scalls} calls.
  (a) Admissions, sessions present, calls per minute? (b) One endpoint: \val{ssvc} s mean service,
  $(\InterarrivalVariationCoefficient^2+\ServiceTimeVariationCoefficient^2)/2 = \val{svar}$. Utilization and mean wait? (c) $\AcceptanceFraction$ falls to
  \val{salphatwo} at the same goodput: new $\ServingUtilization$ and wait?
\item \assumed{} a loop costs \$\val{lturn} per turn; turns one to four accept \val{lqone}, \val{lqtwo},
  \val{lqthree}, \val{lqfour} of tickets reaching them; an engineer resolves the rest for \$\val{lhuman}.
  (a) Where should the loop stop? (b) With \$\val{lhumantwo} triage instead?
\item \assumed{} a child with a \$\val{balloc} allocation and \$\val{bsettled} settled stops locally while a remote
  job bounded at \$\val{bbound} runs. (a) What may the parent release now? (b) The job bills \$\val{blate}: how
  much more returns at close?
\end{qs}
\end{rightcol}
\end{document}
